\label{chap:conclusions}

Este estudo projetou, implementou e avaliou rigorosamente a arquitetura \textbf{TransferKAN} para a classificação automatizada de ferimentos humanos por arma de fogo. Ao combinar a poderosa capacidade de extração de características espaciais de um \textit{backbone} ResNet152 pré-treinado com a modelagem não-linear dinâmica e eficiente em parâmetros de uma Rede de Kolmogorov-Arnold (KAN), estabelecemos um novo e altamente robusto \textit{benchmark} para aplicações de aprendizado profundo em patologia forense.

A análise empírica, conduzida em uma base de dados realista e severamente desbalanceada de 2.551 imagens de cenas de crime (FDCPUnBGunshotDB), demonstrou que o modelo TransferKAN melhora significativamente a robustez e a generalização da classificação em relação aos métodos tradicionais. Especificamente, na tarefa binária de distinguir entre ferimentos de Entrada e de Saída, a arquitetura proposta alcançou uma estelar Área sob a Curva (AUC) micro de 0,928 e uma AUC Macro de 0,908 sob um esquema estrito de validação cruzada de 5 dobras baseado em casos. Esse desempenho superou decisivamente o classificador ResNet152 de referência reportado na literatura de 2024~\cite{Lira2025deep} (que alcançou uma AUC uniforme de 0,821), comprovando que as KANs são amplamente superiores no tratamento probabilístico de classes minoritárias e de fronteiras de características não-lineares, sem colapsar.

Para a classificação da Distância Médico-Legal do Disparo (MLSD), apesar das dificuldades inerentes do desbalanceamento extremo de classes e do protocolo estrito de validação cruzada que previne ativamente o vazamento de dados, o modelo TransferKAN produziu uma AUC Macro (0,749) e uma AUC Ponderada (0,785) mais altas em comparação com as CNNs tradicionais de referência. Esses resultados confirmam fortemente nossa hipótese central: substituir o Perceptron de Múltiplas Camadas (MLP) padrão e rígido por uma Rede de Kolmogorov-Arnold flexível melhora significativamente a capacidade do modelo de discriminar padrões visuais complexos em imagens forenses ruidosas.

Em conclusão, nossa abordagem utilizando a TransferKAN com validação k-fold estrita fornece a metodologia mais cientificamente robusta para a classificação forense de ferimentos por arma de fogo até o momento, demonstrando sua imensa viabilidade como ferramenta de apoio confiável para peritos forenses e médicos-legistas.

Trabalhos futuros devem focar na expansão da base de dados por meio de colaborações multi-institucionais, a fim de validar ainda mais o modelo em diferentes demografias e tipos de arma de fogo. Adicionalmente, uma direção empolgante para pesquisas futuras envolve explorar a interpretabilidade matemática das splines da KAN. Como as funções de ativação são polinômios explícitos por partes nas arestas, elas oferecem uma oportunidade única de desenvolver percepções de IA Explicável (XAI), permitindo potencialmente que patologistas forenses interpretem fisicamente quais características visuais a rede prioriza durante seu processo de decisão. Por fim, investigar a integração de dados multimodais, como a incorporação do texto bruto de laudos de necropsia juntamente com as imagens digitais, poderia refinar ainda mais o \textit{pipeline} de classificação.
